\section{The Agent on Each Lab PC}

\subsection{Standard library only}

\texttt{agent/} imports nothing outside Python's standard library: HTTP with
\texttt{urllib.request}, JSON with \texttt{json}, processes with
\texttt{subprocess}, threads with \texttt{threading}. Installing it is copying
files; updating it is copying newer files. On a fleet of sixty PCs, ``no
\texttt{pip install}'' removes a whole class of failures (no network to PyPI, a
mirror down, a dependency that changed).

\subsection{The main loop}

\begin{lstlisting}[language=Python]
def run_forever(self):
    self.register()
    while True:
        status, hb = self.heartbeat()
        update = hb.get("agent_update")
        if update and self.maybe_update(update):
            sys.exit(updater.EXIT_UPDATED)   # launcher restarts new code
        if not self.healthy or self.busy or hb["command"] in {"drain",
                                                              "quarantine"}:
            time.sleep(5)
            continue
        status, payload = self.claim()       # long poll
        if payload.get("attempt"):
            self.execute(payload["attempt"])
        elif payload.get("build"):
            self.execute_build(payload["build"])
\end{lstlisting}

While a job runs, a second thread renews the lease every few seconds. If a
renew response says \texttt{desired\_state: cancel} (the student pressed Cancel,
or the server's runtime limit passed), the agent sets a \texttt{stop} event and
the executor stops the container.

\subsection{A real bug: a blocking read}

The executor first read job output like this:

\begin{lstlisting}[language=Python]
# before
while True:
    if stop.is_set():
        kill_container(); return 143, "cancelled"
    line = proc.stdout.readline()      # waits here until a line arrives
    if line:
        on_line(line)
    elif proc.poll() is not None:
        break
    if time.time() > deadline:
        kill_container(); return 124, "timeout"
\end{lstlisting}

\texttt{readline()} \emph{blocks}: it does not return until the process prints
another line or exits. A training script that prints ``starting'' and then
trains silently for an hour leaves the loop stuck inside \texttt{readline()}
for that hour. During that time it never checks \texttt{stop} or the deadline,
so \textbf{Cancel did nothing and time limits were ignored}. The unit tests
used chatty scripts and never noticed; the simulated failure exam, which runs a
silent \texttt{time.sleep(120)} job, caught it.

The fix moves reading to a background thread that feeds a queue. The main loop
waits on the queue for at most half a second, so it checks \texttt{stop} and the
deadline twice a second no matter what the job prints:

\begin{lstlisting}[language=Python]
# after (agent/executor.py)
def _pump(proc, on_line, stop, deadline) -> str:
    lines = queue.Queue()

    def reader():
        for line in proc.stdout:
            lines.put(line.rstrip("\n"))
        lines.put(None)                 # end of output

    threading.Thread(target=reader, daemon=True).start()
    while True:
        if stop.is_set():
            return "stop"
        if time.time() > deadline:
            return "timeout"
        try:
            line = lines.get(timeout=0.5)
        except queue.Empty:
            continue
        if line is None:
            proc.wait()
            return "exited"
        on_line(line)
\end{lstlisting}

\paragraph{The general lesson.} Any loop that must react to more than one thing
(a cancel flag, a timer, incoming data) must never block on just one of them.
Other ways to solve the same problem: \texttt{select()} or \texttt{poll()} on
file descriptors, \texttt{asyncio} with timeouts, or Go's \texttt{select} over
channels. The regression tests start a script that prints once and sleeps for
60 seconds, then assert that a 1-second timeout and a cancel both return within
a few seconds. On the old code they hang.

\subsection{Knobs for simulation}

Three environment variables let a GPU-less container pretend to be a lab PC:

\begin{description}
\item[\texttt{LAB\_FAKE\_GPU=ampere-8g:8192:8.6}] report this class, VRAM and
  compute capability. Jobs still run on CPU because \verb|--gpus| is only
  added when \texttt{nvidia-smi} works.
\item[\texttt{LAB\_FAKE\_BUSY\_FILE=/sim/busy}] while the file exists, report
  the desktop GPU as busy.
\item[\texttt{LAB\_FAKE\_DISK\_FREE\_FILE=/sim/disk\_free}] report the number
  in the file as free disk space.
\end{description}

Test seams like these, small, explicit and off by default, let the failure
exam toggle conditions with \texttt{docker exec sim-pc-3 touch /sim/busy}
instead of really filling a disk.
